\begin{lemma}
\label{lemma-extension-normal-domains-and-roots}
Let $A \to B$ be a flat local homomorphism of Noetherian local normal domains.
Let $f \in A$ and $h \in B$ such that $f = w h^n$ for some $n > 1$ and some
unit $w$ of $B$. Assume that for every height $1$ prime
$\mathfrak p \subset A$ there is a height $1$ prime
$\mathfrak q \subset B$ lying over $\mathfrak p$
such that the extension $A_\mathfrak p \subset B_\mathfrak q$ is
weakly unramified. Then $f = u g^n$ for some $g \in A$ and unit $u$ of $A$.
\end{lemma}

\begin{proof}
The local rings of $A$ and $B$ at height $1$ primes are
discrete valuation rings (Algebra, Lemma \ref{algebra-lemma-characterize-dvr}).
Thus the assumption makes sense (via
Definition \ref{definition-extension-discrete-valuation-rings}).
Let $\mathfrak p_1, \ldots, \mathfrak p_r$ be the primes of $A$ minimal
over $f$. These have height $1$ by
Algebra, Lemma \ref{algebra-lemma-minimal-over-1}.
For each $i$ let $\mathfrak q_{i, j} \subset B$, $j = 1, \ldots, r_i$
be the height $1$ primes of $B$ lying over $\mathfrak p_i$.
Say we number them so that  $A_{\mathfrak p_i} \to B_{\mathfrak q_{i, 1}}$
is weakly unramified.
Since $f$ maps to an $n$th power times a unit in $B_{\mathfrak q_{i, 1}}$
we see that the valuation $v_i$ of $f$ in $A_{\mathfrak p_i}$ is
divisible by $n$. Say $v_i = n w_i$ for some $w_i \geq 0$.
Consider the exact sequence
$$
0 \to I \to A \to
\prod\nolimits_{i = 1, \ldots, r}
A_{\mathfrak p_i}/\mathfrak p_i^{w_i}A_{\mathfrak p_i}
$$
defining the ideal $I$. Applying the exact functor $- \otimes_A B$ we obtain
an exact sequence
$$
0 \to I \otimes_A B \to B \to
\prod\nolimits_{i = 1, \ldots, r}
(A_{\mathfrak p_i}/\mathfrak p_i^{w_i}A_{\mathfrak p_i}) \otimes_A B
$$
Fix $i$. We claim that the canonical map
$$
(A_{\mathfrak p_i}/\mathfrak p_i^{w_i}A_{\mathfrak p_i}) \otimes_A B
\to
\prod\nolimits_{j = 1, \ldots, r_i}
B_{\mathfrak q_{i, j}}/\mathfrak q_{i, j}^{e_{i, j}w_i}B_{\mathfrak q_{i, j}}
$$
is injective. Here $e_{i, j}$ is the ramification index of 
$A_{\mathfrak p_i} \to B_{\mathfrak q_{i, j}}$.
The claim asserts that $\mathfrak p_i^{w_i}B_{\mathfrak p_i}$
is equal to the set of elements $b$ of $B_{\mathfrak p_i}$
whose valuation at $\mathfrak q_{i, j}$ is $\geq e_{i, j}w_i$.
Choose a generator $a \in A_{\mathfrak p_i}$ of the principal
ideal $\mathfrak p_i^{w_i}$. Then the valuation of $a$ at
$\mathfrak q_{i, j}$ is equal to $e_{i, j}w_i$. Hence,
as $B_{\mathfrak p_i}$ is a normal domain whose height one
primes are the primes $\mathfrak q_{i, j}$, $j = 1, \ldots, r_i$,
we see that, for $b$ as above, we have $b/a \in B_{\mathfrak p_i}$ by
Algebra, Lemma
\ref{algebra-lemma-normal-domain-intersection-localizations-height-1}.
Thus the claim.

\medskip\noindent
The claim combined with the second exact sequence above
determines an exact sequence
$$
0 \to I \otimes_A B \to B \to
\prod\nolimits_{i = 1, \ldots, r}
\prod\nolimits_{j = 1, \ldots, r_i}
B_{\mathfrak q_{i, j}}/\mathfrak q_{i, j}^{e_{i, j}w_i}B_{\mathfrak q_{i, j}}
$$
It follows that $I \otimes_A B$ is the set of elements $h'$ of $B$
which have valuation $\geq e_{i, j}w_i$ at $\mathfrak q_{i, j}$.
Since $f = wh^n$ in $B$ we see that $h$ has valuation
$e_{i, j}w_i$ at $\mathfrak q_{i, j}$. Thus $h'/h \in B$
by Algebra, Lemma
\ref{algebra-lemma-normal-domain-intersection-localizations-height-1}.
It follows that $I \otimes_A B$ is a free $B$-module of rank $1$
(generated by $h$).
Therefore $I$ is a free $A$-module of rank $1$, see
Algebra, Lemma \ref{algebra-lemma-finite-projective-descends}.
Let $g \in I$ be a generator. Then we see that
$g$ and $h$ differ by a unit in $B$. Working backwards we
conclude that the valuation of $g$ in $A_{\mathfrak p_i}$ is
$w_i = v_i/n$. Hence $g^n$ and $f$ differ by a unit in $A$
(by Algebra, Lemma
\ref{algebra-lemma-normal-domain-intersection-localizations-height-1})
as desired.
\end{proof}